\documentclass[11pt,a4paper]{article}
\usepackage[utf8]{inputenc}
\usepackage[english]{babel}
\usepackage{amsmath}
\usepackage{amssymb}
\usepackage{graphicx}
\usepackage{listings}
\usepackage{xcolor}
\usepackage{hyperref}
\usepackage{geometry}
\usepackage{algorithm}
\usepackage{algpseudocode}

\geometry{margin=2.5cm}

% Python code styling
\lstdefinestyle{pythonstyle}{
    language=Python,
    basicstyle=\ttfamily\small,
    keywordstyle=\color{blue}\bfseries,
    commentstyle=\color{gray}\itshape,
    stringstyle=\color{red},
    showstringspaces=false,
    breaklines=true,
    frame=single,
    numbers=left,
    numberstyle=\tiny\color{gray},
    backgroundcolor=\color{gray!10}
}

\title{Simple Logistic Regression Model for Fake News Detection\\
\large Technical Documentation}
\author{Part 2, Task 2}
\date{\today}

\begin{document}

\maketitle

\section{Overview}

This document provides a detailed explanation of the \texttt{Simple\_model.py} implementation, which creates a baseline logistic regression classifier for binary fake news detection. The model uses a bag-of-words representation with the 10,000 most frequent words as features.

\section{Model Architecture}

\subsection{Binary Classification Task}

The model performs binary classification where:
\begin{itemize}
    \item \textbf{Class 0 (FAKE)}: \texttt{unreliable}, \texttt{hate}, \texttt{junksci}, \texttt{fake}, \texttt{satire}, \texttt{conspiracy}, \texttt{bias}
    \item \textbf{Class 1 (TRUE)}: \texttt{reliable}, \texttt{political}, \texttt{state}, \texttt{clickbait}
\end{itemize}

Let $y \in \{0, 1\}$ denote the binary label for each article.

\subsection{Feature Representation}

The model uses a \textbf{bag-of-words (BoW)} representation:

\begin{equation}
\mathbf{x}_i = [c_1, c_2, \ldots, c_{10000}]^\top \in \mathbb{R}^{10000}
\end{equation}

where $c_j$ is the count of the $j$-th most frequent word in document $i$. This creates a feature matrix:

\begin{equation}
\mathbf{X} \in \mathbb{R}^{n \times 10000}
\end{equation}

where $n$ is the number of documents in the dataset.

\subsection{Logistic Regression Model}

The logistic regression model predicts the probability of an article being TRUE (class 1):

\begin{equation}
P(y=1 | \mathbf{x}) = \sigma(\mathbf{w}^\top \mathbf{x} + b) = \frac{1}{1 + e^{-(\mathbf{w}^\top \mathbf{x} + b)}}
\end{equation}

where:
\begin{itemize}
    \item $\mathbf{w} \in \mathbb{R}^{10000}$ is the weight vector
    \item $b \in \mathbb{R}$ is the bias term
    \item $\sigma(\cdot)$ is the sigmoid function
\end{itemize}

The model is trained by minimizing the binary cross-entropy loss:

\begin{equation}
\mathcal{L}(\mathbf{w}, b) = -\frac{1}{n}\sum_{i=1}^{n} \left[ y_i \log(\hat{y}_i) + (1-y_i)\log(1-\hat{y}_i) \right]
\end{equation}

where $\hat{y}_i = P(y=1 | \mathbf{x}_i)$.

\section{Implementation Details}

\subsection{Pipeline Overview}

The implementation follows a four-stage pipeline:

\begin{algorithm}
\caption{Fake News Detection Pipeline}
\begin{algorithmic}[1]
\State Load train, validation, and test datasets
\State Build vocabulary $V$ of top 10,000 words from training data
\State Create feature matrices $\mathbf{X}_{\text{train}}, \mathbf{X}_{\text{val}}, \mathbf{X}_{\text{test}}$
\State Train logistic regression model on $(\mathbf{X}_{\text{train}}, \mathbf{y}_{\text{train}})$
\State Evaluate on validation and test sets
\State Save trained model and vocabulary
\end{algorithmic}
\end{algorithm}

\subsection{Function Breakdown}

\subsubsection{1. Data Loading: \texttt{load\_data(split)}}

\textbf{Purpose:} Load preprocessed data and create binary labels.

\textbf{Input:} 
\begin{itemize}
    \item \texttt{split}: String indicating which split to load (\texttt{'train'}, \texttt{'val'}, or \texttt{'test'})
\end{itemize}

\textbf{Output:}
\begin{itemize}
    \item DataFrame with columns: \texttt{content} (preprocessed tokens) and \texttt{label} (0 or 1)
\end{itemize}

\textbf{Process:}
\begin{enumerate}
    \item Read CSV file containing preprocessed articles
    \item Remove rows with missing \texttt{content} or \texttt{type}
    \item Map article types to binary labels using the grouping scheme
    \item Filter out articles that don't belong to either class
\end{enumerate}

\begin{lstlisting}[style=pythonstyle, caption={Data Loading Function}]
def load_data(split):
    df = pd.read_csv(f"{DATA_DIR}/{split}.csv", 
                     usecols=['content', 'type'])
    df = df.dropna(subset=['content', 'type'])
    
    df['label'] = df['type'].apply(
        lambda x: 0 if x in FAKE_LABELS 
                 else (1 if x in TRUE_LABELS else -1)
    )
    df = df[df['label'] != -1]
    
    return df[['content', 'label']]
\end{lstlisting}

\subsubsection{2. Vocabulary Building: \texttt{build\_vocabulary(train\_df)}}

\textbf{Purpose:} Extract the 10,000 most frequent words from training data.

\textbf{Input:}
\begin{itemize}
    \item \texttt{train\_df}: Training dataset DataFrame
    \item \texttt{load\_if\_exists}: Boolean flag to load cached vocabulary if available
\end{itemize}

\textbf{Output:}
\begin{itemize}
    \item \texttt{vocab}: Set of 10,000 most frequent words
    \item \texttt{word\_counter}: Counter object with word frequencies (or None if loaded from cache)
\end{itemize}

\textbf{Process:}
\begin{enumerate}
    \item Check if pre-computed vocabulary exists; if so, load and return it
    \item Otherwise, iterate through all training documents
    \item Parse token lists from string representation
    \item Count word frequencies using a Counter
    \item Select top 10,000 most frequent words
    \item Save vocabulary to disk for future use
\end{enumerate}

\textbf{Mathematical Formulation:}

Let $D_{\text{train}}$ be the set of training documents, and $w_{ij}$ be the $j$-th word in document $i$. The word frequency is:

\begin{equation}
f(w) = \sum_{i=1}^{|D_{\text{train}}|} \sum_{j=1}^{|d_i|} \mathbb{1}[w_{ij} = w]
\end{equation}

The vocabulary is then:

\begin{equation}
V = \{w_1, w_2, \ldots, w_{10000}\} \text{ where } f(w_1) \geq f(w_2) \geq \cdots \geq f(w_{10000})
\end{equation}

\begin{lstlisting}[style=pythonstyle, caption={Vocabulary Building Function}]
def build_vocabulary(train_df, load_if_exists=True):
    vocab_file = f"{DATA_DIR}/top_{TOP_K_WORDS}_vocab.pkl"
    
    if load_if_exists and pd.io.common.file_exists(vocab_file):
        with open(vocab_file, 'rb') as f:
            return pickle.load(f), None
    
    word_counter = Counter()
    for content in train_df['content']:
        try:
            tokens = ast.literal_eval(content)
            word_counter.update(tokens)
        except:
            continue
    
    vocab = set([word for word, _ 
                 in word_counter.most_common(TOP_K_WORDS)])
    
    with open(vocab_file, 'wb') as f:
        pickle.dump(vocab, f)
    
    return vocab, word_counter
\end{lstlisting}

\subsubsection{3. Feature Creation: \texttt{create\_features(df, vocab)}}

\textbf{Purpose:} Convert documents into bag-of-words feature vectors.

\textbf{Input:}
\begin{itemize}
    \item \texttt{df}: DataFrame containing documents
    \item \texttt{vocab}: Set of vocabulary words
\end{itemize}

\textbf{Output:}
\begin{itemize}
    \item \texttt{X}: Feature matrix of shape $(n, 10000)$ where $n$ is the number of documents
\end{itemize}

\textbf{Process:}
\begin{enumerate}
    \item Create a mapping from words to column indices
    \item Initialize a zero matrix of shape $(n, 10000)$
    \item For each document:
    \begin{itemize}
        \item Parse the token list
        \item Count occurrences of each word
        \item For words in vocabulary, set the corresponding column to the count
    \end{itemize}
\end{enumerate}

\textbf{Mathematical Formulation:}

For document $i$ with tokens $\{w_{i1}, w_{i2}, \ldots, w_{im}\}$, the feature vector is:

\begin{equation}
x_{ij} = \sum_{k=1}^{m} \mathbb{1}[w_{ik} = v_j]
\end{equation}

where $v_j$ is the $j$-th word in the vocabulary $V$.

\begin{lstlisting}[style=pythonstyle, caption={Feature Creation Function}]
def create_features(df, vocab):
    vocab_to_idx = {word: idx 
                    for idx, word in enumerate(sorted(vocab))}
    X = np.zeros((len(df), len(vocab)), dtype=np.int32)
    
    for i, content in enumerate(df['content']):
        try:
            tokens = ast.literal_eval(content)
            for word, count in Counter(tokens).items():
                if word in vocab_to_idx:
                    X[i, vocab_to_idx[word]] = count
        except:
            continue
    
    return X
\end{lstlisting}

\subsubsection{4. Main Pipeline: \texttt{main()}}

\textbf{Purpose:} Orchestrate the entire training and evaluation pipeline.

\textbf{Process:}
\begin{enumerate}
    \item \textbf{Load Data:} Load train, validation, and test splits
    \item \textbf{Build Vocabulary:} Extract top 10,000 words from training data
    \item \textbf{Create Features:} Convert all splits to feature matrices
    \item \textbf{Train Model:} Fit logistic regression on training data
    \item \textbf{Validate:} Evaluate on validation set
    \item \textbf{Test:} Evaluate on test set (final performance)
    \item \textbf{Save:} Persist trained model and vocabulary
\end{enumerate}

\begin{lstlisting}[style=pythonstyle, caption={Main Pipeline}]
def main():
    # Load data
    train_df = load_data('train')
    val_df = load_data('val')
    test_df = load_data('test')
    
    # Build vocabulary from training data only
    vocab, word_counter = build_vocabulary(train_df)
    
    # Create feature matrices
    X_train = create_features(train_df, vocab)
    y_train = train_df['label'].values
    
    X_val = create_features(val_df, vocab)
    y_val = val_df['label'].values
    
    X_test = create_features(test_df, vocab)
    y_test = test_df['label'].values
    
    # Train logistic regression
    model = LogisticRegression(max_iter=1000, random_state=42)
    model.fit(X_train, y_train)
    
    # Evaluate on validation set
    y_val_pred = model.predict(X_val)
    print(f"Val F1: {f1_score(y_val, y_val_pred):.4f}")
    
    # Evaluate on test set
    y_test_pred = model.predict(X_test)
    print(f"Test F1: {f1_score(y_test, y_test_pred):.4f}")
    
    # Save model
    with open(f"{DATA_DIR}/logistic_model.pkl", 'wb') as f:
        pickle.dump(model, f)
\end{lstlisting}

\section{Hyperparameters}

\begin{table}[h]
\centering
\begin{tabular}{|l|l|p{7cm}|}
\hline
\textbf{Parameter} & \textbf{Value} & \textbf{Description} \\
\hline
Vocabulary Size & 10,000 & Number of most frequent words used as features \\
\hline
Solver & lbfgs & Limited-memory BFGS optimization algorithm \\
\hline
Max Iterations & 1,000 & Maximum number of iterations for convergence \\
\hline
Random State & 42 & Seed for reproducibility \\
\hline
Feature Type & Word Counts & Raw counts (no TF-IDF weighting) \\
\hline
\end{tabular}
\caption{Model Hyperparameters}
\end{table}

\section{Data Splits}

The dataset is split using a deterministic hash-based approach:

\begin{itemize}
    \item \textbf{Training Set:} 80\% of data ($\sim$800,000 samples)
    \item \textbf{Validation Set:} 10\% of data ($\sim$100,000 samples)
    \item \textbf{Test Set:} 10\% of data ($\sim$100,000 samples)
\end{itemize}

The splits are stratified to maintain class balance across all sets.

\section{Evaluation Metrics}

The model is evaluated using:

\begin{itemize}
    \item \textbf{F1 Score:} Harmonic mean of precision and recall
    \begin{equation}
    F1 = 2 \cdot \frac{\text{Precision} \cdot \text{Recall}}{\text{Precision} + \text{Recall}}
    \end{equation}
    
    \item \textbf{Accuracy:} Proportion of correct predictions
    \begin{equation}
    \text{Accuracy} = \frac{\text{TP} + \text{TN}}{\text{TP} + \text{TN} + \text{FP} + \text{FN}}
    \end{equation}
    
    \item \textbf{Classification Report:} Per-class precision, recall, and F1 scores
\end{itemize}

where TP = True Positives, TN = True Negatives, FP = False Positives, FN = False Negatives.

\section{Key Design Decisions}

\subsection{Why Bag-of-Words?}

\begin{itemize}
    \item \textbf{Simplicity:} Easy to implement and interpret
    \item \textbf{Baseline:} Provides a simple baseline for comparison with advanced models
    \item \textbf{Efficiency:} Fast to compute and train
\end{itemize}

\subsection{Why 10,000 Words?}

\begin{itemize}
    \item \textbf{Coverage:} Captures most informative words while avoiding noise
    \item \textbf{Computational Efficiency:} Keeps feature space manageable
    \item \textbf{Generalization:} Reduces overfitting compared to using all words
\end{itemize}

\subsection{Why No TF-IDF?}

The assignment explicitly states: ``You do not need to apply TF-IDF weighting to the features.'' This keeps the baseline simple and allows for easier comparison when TF-IDF is added later.

\subsection{Why Training Data Only for Vocabulary?}

Building vocabulary from training data only prevents \textbf{data leakage}. If we used validation or test data, the model would have indirect access to information from those sets, leading to overly optimistic performance estimates.

\section{Limitations and Future Improvements}

\subsection{Current Limitations}

\begin{enumerate}
    \item \textbf{No Word Order:} Bag-of-words ignores word sequence and context
    \item \textbf{No Semantic Understanding:} Words are treated as independent features
    \item \textbf{Sparse Representation:} Most feature values are zero (inefficient)
    \item \textbf{Fixed Vocabulary:} Cannot handle new words at test time
\end{enumerate}

\subsection{Potential Improvements}

\begin{enumerate}
    \item \textbf{TF-IDF Weighting:} Down-weight common words, up-weight rare informative words
    \item \textbf{N-grams:} Capture local word order (bigrams, trigrams)
    \item \textbf{Word Embeddings:} Use pre-trained embeddings (Word2Vec, GloVe)
    \item \textbf{Deep Learning:} LSTM, CNN, or Transformer-based models
    \item \textbf{Meta-data Features:} Include domain, author, publication date (Part 2, Task 3)
\end{enumerate}

\section{Expected Performance}

According to the assignment guidelines:
\begin{itemize}
    \item \textbf{Training Time:} Should take no more than 5 minutes on a modern laptop
    \item \textbf{Expected F1 Score:} Approximately 94\% on the test set
    \item Note: Actual performance may vary based on data splitting and label grouping decisions
\end{itemize}

\section{Conclusion}

This implementation provides a clean, minimal baseline for fake news detection using logistic regression with bag-of-words features. The modular design allows for easy extension with additional features, different models, or preprocessing techniques in subsequent tasks.

\end{document}
